\documentclass[12pt]{article}
\usepackage[utf8]{inputenc}
\usepackage[T1]{fontenc}
\usepackage{hyperref}
\title{Apalachee Language Resources for K-12 Curriculum Development}
\author{Compiled by Claude Sonnet 5.5 (claude-sonnet-5-5) for Andr\'as Kornai}
\date{}
\begin{document}
\maketitle

\section{Introduction}
Apalachee (ISO 639-3 xap; Glottocode apal1237) is an extinct Muskogean language of the United States. Searches found no living speakers, no school programme and no revitalization curriculum for it. Sources disagree on when it ceased to be spoken (early 18th century to the 19th century, among Apalachee people removed to Louisiana). Very little material exists, so any curriculum would have to be built from a small documentary base and from related living Muskogean languages.

\section{Primary Educational Resources}
No primers, textbooks or classroom materials for Apalachee turned up in the searches. None are listed here.

\section{Linguistic and Reference Materials}
The language is known mainly from a single document, a letter written in 1688 to Charles II of Spain. Geoffrey Kimball produced a grammatical sketch and a vocabulary based on it. The vocabulary is described as giving Spanish equivalents and Muskogean cognates where possible.
\begin{itemize}
\item Wikipedia, ``Apalachee language'': overview of the language and its documentation. \url{https://en.wikipedia.org/wiki/Apalachee_language}
\item ``An Apalachee Vocabulary'' (UBC Wiki, RelLex): a vocabulary page. \url{https://wiki.ubc.ca/Documentation:RelLex/An_Apalachee_Vocabulary}
\item Endangered Languages Project resources page, listed under the node for Apalachee. I could not confirm its contents. \url{https://endangeredlanguages.com/node/38123/resources}
\end{itemize}

\section{Religious and Community Resources}
No Apalachee-language scripture or community language programme was found. Apalachee is said to be closely related to Koasati and Alabama, whose communities have their own programmes. This search did not find those programmes.

\section{Digital and Audio Resources}
No audio recordings of Apalachee exist, since the language is known only from written sources. For comparison, the search returned material on the revival of the related Muscogee (Creek) language, which has an interactive corpus.
\begin{itemize}
\item Muscogee Soundwaves: An Interactive Corpus (a related language, not Apalachee). \url{https://solve.mit.edu/solutions/90157}
\item The Revival of the Muscogee Creek Language through the Development of an Indigenous Language Interest Group (related language, not Apalachee). \url{https://scholars.lynn.edu/en/publications/the-revival-of-the-muscogee-creek-language-through-the-developmen/}
\end{itemize}

\section{Curriculum Development Context}
A K-12 programme cannot use Apalachee as a language of instruction. Realistic options are a heritage-awareness unit with the attested vocabulary, or reference to the revitalization methods used for related Muskogean languages such as Muscogee. Any use of the material should be arranged with descendant Apalachee communities, which these searches did not identify.
\end{document}